\section{MoE 基础：专家、路由器与条件参数}

明确资源边界后，才能定义 Mixture of Experts（MoE，专家混合）。这里的“专家”不是人工指定的语言学角色，而是若干可学习的子网络；“混合”也不意味着每次都平均执行全部子网络，而是由 router 根据输入选择少量路径。这个思想早于 Transformer，1991 年的 Adaptive Mixtures of Local Experts 已经提出由 gating network 学习分工。

\lecturefigure{slide-04-adaptive-mixtures-history.jpg}{从 Adaptive Mixtures of Local Experts 到稀疏神经语言模型的历史线索。}{Stanford Online 官方视频 00:03:06--00:04:01。}

\begin{knowledgebox}{术语表：先把路由语言说清楚}
\begin{tabularx}{\textwidth}{L{0.22\textwidth}X}
\toprule
术语 & 本讲中的精确定义 \\
\midrule
MoE & 一组 expert 加一个 router；输入只激活少量 expert，因此参数条件使用。 \\
expert & 通常是独立的 feed-forward network（FFN）参数块，不等于人工可解释职业。 \\
router / gating & 从 token 表示计算 expert 分数或概率，并决定 dispatch 路径的轻量网络。 \\
top-1 / top-2 & 每个 token 分别选择一个或两个最高分 expert；后者通常增加计算与通信。 \\
sparse activation & 大部分 expert 对当前 token 不执行；它不同于把权重永久剪枝为零。 \\
load balance & 让各 expert 接收的 token 数接近，便于形成规则的批量矩阵乘法。 \\
\bottomrule
\end{tabularx}
\end{knowledgebox}

设 token 表示为 $x\in\mathbb{R}^{d}$，有 $E$ 个 experts，router 参数为 $W_r$。标准 gating 可以写为：

\[
p_i(x)=\operatorname{softmax}(W_rx)_i,
\qquad
y(x)=\sum_{i\in\mathcal{S}(x)} g_i(x)E_i(x).
\]

其中，$p_i(x)$ 是 token 被分给第 $i$ 个 expert 的 router 概率；$\mathcal{S}(x)$ 是 top-$k$ 选出的 expert 集合；$g_i(x)$ 是选中后用于缩放输出的权重；$E_i(x)$ 是第 $i$ 个 expert 的 FFN 输出；$y(x)$ 是稀疏层输出。只有 $\mathcal{S}(x)$ 中的 experts 被执行。

\lecturefigure{slide-05-moe-gating-equations.jpg}{MoE 的 gating 概率、expert 选择与加权输出。}{Stanford Online 官方视频 00:04:02--00:04:25。}

\begin{importantbox}{MoE 的本质}
MoE 不是“把多个完整模型投票平均”，而是对每个 token 运行一个很小的 router，再调用少量 FFN 参数块。它让不同 token 使用不同权重，却可以保持相近的 active compute。
\end{importantbox}

\teachervoice{Irwan Bello 的口头定义很重要：每个输入大致接受相似数量的计算，但被应用的是不同权重。这个表述把 sparsity 与“简单减少计算步骤”区分开，也为后面的 adaptive computation 讨论埋下伏笔。}

早期稀疏 MoE 在机器翻译上取得明显成功，但实际使用需要解决三个系统问题。第一，路由让 token 数量随 expert 动态变化；第二，expert 跨设备放置时会产生通信；第三，router、低精度和负载波动会放大训练不稳定。Switch Transformer 的贡献不是首先提出 MoE，而是把这些障碍压缩成更简单、可训练、可扩展的方案。

\lecturefigure{slide-06-moe-success-and-challenges.jpg}{早期 MoE 的成功及其复杂度、通信与稳定性障碍。}{Stanford Online 官方视频 00:04:26--00:05:37。}

\begin{warningbox}{不要把 expert specialization 当作先验事实}
router 可能形成主题、语言或句法上的分工，也可能只是学习到更难解释的数值分区。专家名字通常是分析结果，不是训练前约束；看到某个 token 更常去某 expert，也不能直接推出该 expert 具有稳定、因果的“技能”。
\end{warningbox}

\subsection{本章小结}

MoE 由 expert 参数块和 router 组成，核心是按 token 条件选择参数。它扩大可用容量，却引入动态形状、通信、稳定性与负载均衡问题。

